\documentclass[9pt,aspectratio=169]{beamer}

% --- Theme ---
\usetheme{Singapore}
\setbeamertemplate{navigation symbols}{}   % no icons
\setbeamertemplate{mini frames}{}          % no dots in headline
\setbeamercovered{transparent}

% --- Grayscale colors ---
\definecolor{PelloGrayDark}{RGB}{60,60,60}
\definecolor{PelloGray}{RGB}{120,120,120}
\definecolor{PelloGrayLight}{RGB}{235,235,235}
\definecolor{PelloGrayUltraLight}{RGB}{248,248,248}

\usecolortheme[named=PelloGrayDark]{structure}

% --- Blocks ---
\setbeamertemplate{blocks}[rounded][shadow=false]
\setbeamercolor{block title}{fg=white,bg=PelloGrayDark}
\setbeamercolor{block body}{fg=black,bg=PelloGrayLight}
\setbeamercolor{block title example}{fg=white,bg=PelloGrayDark}
\setbeamercolor{block body example}{fg=black,bg=PelloGrayLight}
\setbeamercolor{section in head/foot}{fg=PelloGray,bg=PelloGrayLight}
\setbeamercolor{subsection in head/foot}{fg=PelloGray,bg=PelloGrayLight}
\setbeamercolor{frametitle}{fg=black}

% --- Itemize bullets (optional; keep if you like the style) ---
\setbeamertemplate{itemize item}{\raisebox{0.2em}{\scalebox{.6}{$\bullet$}}}
\setbeamertemplate{itemize subitem}{\raisebox{0.2em}{\scalebox{1}{$-$}}}

% --- Footer: frame numbers ---
\makeatletter
\setbeamertemplate{footline}{%
  \hfill\usebeamerfont{footline}\insertframenumber/\inserttotalframenumber\hspace{1em}\vspace{2pt}
}
\makeatother

% --- Minimal packages (add more only when needed) ---
\usepackage[english]{babel}
\usepackage[T1]{fontenc}
\usepackage[utf8]{inputenc}
\usepackage{graphicx}
\usepackage{tikz}
\usepackage{booktabs}
\usepackage{makecell} % for jumping lines within cells in tables.
\usepackage{multirow}
\usepackage{amsmath}
\usepackage{hyperref}
\hypersetup{colorlinks,linkcolor=,urlcolor=blue}
\usepackage{comment}
\usepackage{listings}
\lstset{
  language=R,
  basicstyle=\ttfamily\footnotesize,
  keywordstyle=\color{blue},
  commentstyle=\color{green!50!black},
  stringstyle=\color{red},
  frame=single,
  showstringspaces=false
}
\linespread{1.7}   % 1.0 = normal, 1.2–1.5 common
\AtBeginEnvironment{lstlisting}{\linespread{1}\selectfont} % keep code listings single-spaced

% --- Meta ---
\title{Applied Economics}
\subtitle{Introduction}
\author{Pello Aspuru}
%\institute{University of Vienna}
\date{Winter 2026/2027}

% --- Title page spacing + bottom-left author/date ---
\setbeamertemplate{title page}{
  \vspace{8em}
  \begin{center}
	{\usebeamerfont{title}\inserttitle\par}
  	\vspace{1.5em} % <-- increase title–subtitle gap
  {\usebeamerfont{subtitle}\insertsubtitle\par}
  \end{center}
  \vspace{4em}
  \begin{flushleft}
    {\usebeamerfont{author}\insertauthor\par}
    {\usebeamerfont{date}\insertdate}
  \end{flushleft}
}

\begin{document}

% Title slide without headline/footer
{
  \setbeamertemplate{headline}{}
  \setbeamertemplate{footline}{}
  \begin{frame}
    \titlepage
  \end{frame}
}


\section[Course Logistics]{Introduction}

\begin{frame}[fragile]{Welcome to Applied Economics}

%Lectures
%\begin{itemize}
%    \item Wednesday (040178-1) 16:45-18:15, PC-Seminarraum 5
%    \item Thursday (040178-2)  16:45-18:15, PC-Seminarraum 3 
%\end{itemize}


All course material (including these slides) will be distributed via \textbf{moodle}.

\alert{Prior knowledge at the level of Introductory Econometrics (BA) will be assumed throughout the course.}

%For course structure and deadlines, see the \href{syllabus}{syllabus}

Class participation is highly encouraged, determines 10\% of final grade.

Office hours: by email appointment via pello.aspuru@univie.ac.at

\onslide<2>{
People who signed up online for this class and signed the list today
are considered \textbf{enrolled}.

If you are registered, you will receive a grade!
}
\end{frame}

\begin{frame}{Course Overview}
 \begin{block}{In this Course You Will Learn to...}
\begin{itemize} 
\item Use econometric tools to produce empirical analyses.
\item Learn to interpret, and critically evaluate existing empirical studies.
\item Understand possibilities and limitations of empirical analyses.
\end{itemize}
\end{block}
\end{frame}

\begin{frame}{Course Tools}
\onslide<1->{
\begin{block}{Statistical Software: R and RStudio}
\begin{itemize} 
\item Freely available open source software
\item Active user community providing up-to-date functions \& support
\item Extensive statistical modeling capabilities
\item Relevant for industry jobs
\end{itemize}
\end{block}}

\onslide<2->{
\begin{block}{Auxiliary Software}
\begin{itemize} 
\item Large Language Models (ChatGPT, Claude etc.) to support learning \& coding.
\item Latex for drafting academic-style documents, problem set.
\end{itemize}
\end{block}}

\end{frame}

\begin{frame}{Course Assessment}

\begin{block}{Class Participation (10\%)}
\begin{itemize} 
\item[] Classroom engagement through questions \& comments is encouraged!
\end{itemize}
\end{block}

\begin{block}{Midterm Exam (25\%)}
\begin{itemize}
%\item based on classes, slides, and problem sets
\item[] Preliminary date: November 25th/26th 2026.\\
%\item Preliminary location: TBD.
%\item there will be no resit
\end{itemize}
\end{block}

\onslide<2->{
\begin{block}{Project Presentations (25\%)}
\begin{itemize}
\item Work on the project in groups of up to 3 students, fill in names using the link provided by Moodle.
\item Presentation I (10\%): December 16th/17th 2026.
\item Presentation II (15\%): January 20th/21st 2027.
\end{itemize}
\end{block}

\begin{block}{Final exam (40\%)}
\begin{itemize} 
%\item based on classes, slides, and problem sets
\item[] Date: Wednesday, 27th January 2027. \\
%\item Preliminary location: lecture hall 9, 01.314 Oskar-Morgenstern-Platz 1, OG 1st floor
%\item there will be no resit
\end{itemize}
\end{block}}

%All tasks are compulsory. Missing \textit{any} of them implies failing the course.
\end{frame}

\begin{frame}{Useful Textbooks}
\linespread{1.3}\selectfont % dense slide: tighter line spacing so it fits
  \begin{enumerate}
    \item \textbf{Heiss, F. (2020).} Using R for Introductory Econometrics (2nd edition)\\
    available online: \href{http://urfie.net/readOnline.html}{http://urfie.net/readOnline.html}

    \item \textbf{Hanck, C., Arnold, M., Gerber, A., \& Schmelzer, M. (2023).} Introduction to Econometrics with R\\
    available online: \href{https://www.econometrics-with-r.org/index.html}{https://www.econometrics-with-r.org/index.html}

    \item \textbf{Oswald, F., Viers, V., Villedieu, P., \& Kennedy, G. (2020).}
    Introduction to Econometrics with R. Sciences Po, Department of Economics, Paris, France.\\
    available online: \href{https://scpoecon.github.io/ScPoEconometrics/}{https://scpoecon.github.io/ScPoEconometrics/}.
  \end{enumerate}

  \vspace{0.3cm}
  If you need to read up on theoretical concepts:
  \begin{itemize}
    \item Angrist, J. D., \& Pischke, J.-S. (2008). Mostly Harmless Econometrics. % Princeton University Press.
    \item Wooldridge, J. (2015). Introductory Econometrics: A Modern Approach. % (Sixth edition, student ed.)
    % \item Hill, R. C., Griffiths, W. E., \& Lim, G. C. (2018). Principles of Econometrics (5th ed.)
  \end{itemize}
\end{frame}

\section{Motivation}

\begin{frame}{What is Applied Economics?}
\linespread{1.5}\selectfont % dense slide: tighter line spacing so it fits

\begin{center}
  The use of economic theory and \textbf{data} to answer real-world \textbf{economic} questions.
\end{center}

\begin{itemize}
\onslide<2->{
\item \textcolor{teal}{Test} what economic theory predicts.
\begin{itemize}
\item Theory: a higher minimum wage reduces employment. Do we see this in the data?
\end{itemize}}
\onslide<3->{
\item \textcolor{teal}{Discover} (causal) relationships between economic variables and \alert{quantify} them.
\begin{itemize}
%\item Does more education \textcolor{teal}{increase earnings}? \alert{If so, by how much?}
\item Does immigration \textcolor{teal}{increase unemployment} for locals? \alert{By how much?}
\end{itemize}}
\onslide<4->{
\item \textcolor{teal}{Evaluate and design} policies that achieve the intended objective at the \textbf{lowest possible cost.}
\begin{itemize}
%\item Do job training programmes for the unemployed affect the chance to find a new job?
\item How much money do we need to subsidize public transport to save 1 ton of carbon emissions?
\end{itemize}}
\end{itemize}

\onslide<5->{
\begin{block}{}
  \begin{enumerate}
    \item Gather, clean and combine multiple sources of data $\rightarrow$ Tidy dataset.
    \item Find patterns in the data $\rightarrow$ Data Visualization and Descriptive Statistics.
    \item Apply econometric tools $\rightarrow$ Establish causality and estimate magnitude.
  \end{enumerate}
\end{block}

}

\end{frame}




%\begin{frame}{What We Will Focus on in This Course}
%How to estimate coefficients of the MLR using R?

%When do the coefficients of the MLR reflect actual causal effects of interest?

%How to do statistical inference?

%What are frequently used modern \textit{quasi-experimental methods} and how can we apply them in R? 

%\end{frame}

\begin{comment}
\section{Course Outline}

\begin{frame}[fragile]{Course Outline}
1. Introduction \textit{(Heiss, Sections 1 - 1.4, 1.6)}
 \begin{itemize}
   \item Motivation
    \item Introduction to R
\end{itemize}

2. Revisiting the Multiple Linear Regression Model \\
\textit{(Heiss, 3.1 and 3.3; Wooldridge 3.1--3.2)}
\begin{columns}
    \begin{column}{0.38\linewidth}
    \begin{itemize}
     \item Structure and terminology 
    \item Use for causal analysis  
    \end{itemize} 
    \end{column}
        \begin{column}{0.46\linewidth}
\begin{itemize}
   
    \item Matrix Notation
    \item Deriving the OLS estimator
\end{itemize} 
    \end{column}
\end{columns}

3. Revisiting the Multiple Linear Regression Model II \\
\textit{(Heiss, 6.1 and 7; Wooldridge 3.3 and 5.1-5.2)}
\begin{columns}
    \begin{column}{0.38\linewidth}
    \begin{itemize}
    \item MLR Assumptions 
    \item Properties of the OLS estimator
     \end{itemize} 
    \end{column}
        \begin{column}{0.46\linewidth}
\begin{itemize}
    \item Accommodating nonlinear relationships
\end{itemize}
    \end{column}
\end{columns}
%12 lectures
%PSET due dates
\end{frame}

\begin{frame}[fragile]{Course Outline}
4. - OLS Variance \textit{(Heiss, 3.4, 8.1; Hanck et al. 5.4)} \\ \hspace*{+.27cm} - more R topics: \texttt{ggplot2}, \texttt{for-loops} \textit{(Heiss, 1.9)}

5. Inference I \textit{(Heiss, 4.1-4.2; Wooldridge, Ch. 4.1-4.3)}
    \begin{itemize}
     \item t-Tests, p-values, confidence intervals
    \end{itemize} 
    
6. - Inference II: \textit{(Heiss, 4.3-4.4, 8.1; Wooldridge, Ch. 4.4-4.5)}
\vspace{-0.05cm}
 \begin{itemize}
     \item Multiple Hypothesis Testing
    \end{itemize} 
    \vspace{+0.15cm}

\\ \hspace*{+.275cm} - More R topics: writing own functions    \textit{(Heiss, 1.9.3)}
   
%12 lectures
%PSET due dates
\end{frame}

\begin{frame}[fragile]{Course Outline}
7. Endogeneity  \textit{(Hanck et al. 9-9.2, 9.4)}

8. Instrumental Variable Estimation \textit{(Oswald et al. Ch. 11)}

9. Panel Data methods \textit{(Oswald et al. Ch. 12)}

10. Diff-in-Diff Estimator \textit{(Hanck et al. Ch. 13.4)}

11. Regression Discontinuity Design \textit{(Hanck et al. Ch. 13.4)}
%12 lectures
%PSET due dates
\end{frame}
\end{comment}



\end{document}
